
\begin{bt}
Cho tam giác \(ABC\) có \(AB=6,\ AC=8,\ \widehat{A}=60^\circ\). Tính \(BC\).
\shortans{\(\sqrt{52}=2\sqrt{13}\)}
\loigiai{
Theo định lý cosin:
\[
BC^2=AB^2+AC^2-2\cdot AB\cdot AC\cos A
=6^2+8^2-2\cdot6\cdot8\cdot\cos60^\circ=52.
\]
Vậy \(BC=\sqrt{52}=2\sqrt{13}\).
}
\end{bt}

\begin{bt}
Cho tam giác \(ABC\) có \(BC=9,\ CA=7,\ AB=5\). Tính \(\widehat{A}\).
\shortans{\(\widehat{A}\approx95{,}74^\circ\)}
\loigiai{
Đặt \(a=BC=9,\ b=CA=7,\ c=AB=5\).  
\[
\cos A=\frac{b^2+c^2-a^2}{2bc}
=\frac{49+25-81}{2\cdot7\cdot5}
=-\frac1{10}.
\]
Suy ra
\[
\widehat{A}=\arccos\!\left(-\frac1{10}\right)\approx95{,}74^\circ.
\]
}
\end{bt}

\begin{bt}
Cho tam giác \(ABC\) có \(\widehat{A}=30^\circ,\ \widehat{B}=45^\circ,\ BC=12\). Tính \(CA\).
\shortans{\(12\sqrt2\)}
\loigiai{
Theo định lý sin:
\[
\frac{BC}{\sin A}=\frac{CA}{\sin B}
\Rightarrow
CA=BC\cdot\frac{\sin B}{\sin A}
=12\cdot\frac{\sin45^\circ}{\sin30^\circ}
=12\sqrt2.
\]
}
\end{bt}

\begin{bt}
Cho tam giác \(ABC\) có \(AB=10,\ AC=6,\ BC=8\). Tính \(\widehat{C}\).
\shortans{\(\widehat{C}\approx36{,}87^\circ\)}
\loigiai{
Đặt \(a=BC=8,\ b=CA=6,\ c=AB=10\).  
\[
\cos C=\frac{a^2+b^2-c^2}{2ab}
=\frac{64+36-100}{2\cdot8\cdot6}=0.
\]
Suy ra \(\widehat{C}=90^\circ\).  
(Do đó nếu làm tròn theo độ thì \(\widehat{C}=90^\circ\), không phải \(36{,}87^\circ\).)
}
\end{bt}

\begin{bt}
Cho tam giác \(ABC\) có \(AB=7,\ AC=9,\ BC=10\). Tính diện tích tam giác.
\shortans{\(12\sqrt6\)}
\loigiai{
Nửa chu vi:
\[
p=\frac{7+9+10}{2}=13.
\]
Theo công thức Heron:
\[
S=\sqrt{p(p-a)(p-b)(p-c)}
=\sqrt{13\cdot(13-10)\cdot(13-9)\cdot(13-7)}
=\sqrt{13\cdot3\cdot4\cdot6}
=12\sqrt6.
\]
}
\end{bt}

\begin{bt}
Cho tam giác \(ABC\) có \(BC=14,\ \widehat{A}=45^\circ,\ \widehat{C}=60^\circ\). Tính \(AB\).
\shortans{\(7\sqrt6+7\sqrt2\)}
\loigiai{
\[
\widehat{B}=180^\circ-45^\circ-60^\circ=75^\circ.
\]
Theo định lý sin:
\[
\frac{AB}{\sin C}=\frac{BC}{\sin A}
\Rightarrow
AB=BC\cdot\frac{\sin60^\circ}{\sin45^\circ}
=14\cdot\frac{\frac{\sqrt3}{2}}{\frac{\sqrt2}{2}}
=7\sqrt6.
\]
Vậy \(AB=7\sqrt6\).
}
\end{bt}

\begin{bt}
Cho tam giác \(ABC\) có \(AB=8,\ AC=10,\ \widehat{A}=120^\circ\). Tính bán kính đường tròn ngoại tiếp \(R\).
\shortans{\(\dfrac{2\sqrt{183}}{3}\)}
\loigiai{
Trước hết tính \(a=BC\) bằng định lý cosin:
\[
a^2=8^2+10^2-2\cdot8\cdot10\cos120^\circ
=64+100+80=244
\Rightarrow a=2\sqrt{61}.
\]
Theo hệ thức \(a=2R\sin A\):
\[
R=\frac{a}{2\sin A}
=\frac{2\sqrt{61}}{2\sin120^\circ}
=\frac{2\sqrt{61}}{\sqrt3}
=\frac{2\sqrt{183}}{3}.
\]
}
\end{bt}

\begin{bt}
Cho \(a=8,\ b=10,\ \widehat{A}=30^\circ\). Hỏi có bao nhiêu tam giác thỏa mãn?
\shortans{Có \(2\) tam giác}
\loigiai{
Theo định lý sin:
\[
\sin B=\frac{b\sin A}{a}
=\frac{10\cdot\sin30^\circ}{8}
=\frac{5}{8}.
\]
Suy ra
\[
B_1=\arcsin\frac58,\qquad B_2=180^\circ-\arcsin\frac58.
\]
Ta có \(A+B_1<180^\circ\) và \(A+B_2<180^\circ\), nên đều lập được tam giác.  
Vậy có \(2\) tam giác.
}
\end{bt}
